\documentclass[11pt]{article}
\usepackage[utf8]{inputenc}
\usepackage[T1]{fontenc}
\usepackage{amsmath,amssymb,amsthm}
\usepackage{booktabs}
\usepackage{array}
\usepackage[margin=1.2in]{geometry}
\usepackage{url}

\newtheorem{theorem}{Theorem}
\newtheorem{proposition}[theorem]{Proposition}
\newtheorem{corollary}[theorem]{Corollary}
\newtheorem{lemma}[theorem]{Lemma}
\newtheorem{conjecture}[theorem]{Conjecture}
\newtheorem{definition}[theorem]{Definition}
\newtheorem{remark}[theorem]{Remark}


\title{The smallest semi-bimagic square of order six}
\author{Luca Cipolletta}
\date{\today}

\begin{document}
\maketitle

\begin{abstract}
A $6\times 6$ array of distinct positive integers is \emph{semi-bimagic}
if all six rows and all six columns have the same sum $S_1$ and the same
sum of squares $S_2$ (diagonals unconstrained). The oldest example, due
to Pfeffermann (1894), uses integers up to $49$. We show by exhaustive
computer enumeration that the smallest possible largest entry is
$\mathrm{MaxNb}=42$: semi-bimagic squares with distinct integers in
$[1,42]$ exist, and none exists with all entries in $[1,41]$. At the
optimum, $S_1=129$ is forced, and exactly four $36$-element sets of
integers admit such a square, with $S_2\in\{3659,3671,3691,3799\}$; all
four sets are closed under $x\mapsto 43-x$, a symmetry that emerges from
the enumeration rather than being imposed. All exclusions $36$--$41$
were confirmed by two independent algorithms with different underlying
assumptions, and the classification at the optimum by two independent
implementations jointly covering the full admissible window
$S_1\in[112,147]$. As a corollary, squaring the minimal square cell-wise
yields a $6\times6$ semi-magic square of $36$ distinct perfect squares
with constant $3659$.

A second, structural part isolates the local mechanism by which two
columns of a semi-bimagic square interact: \emph{twin triples},
three-element subsets of the two column sets with equal sums and equal
sums of squares (degree-$2$ Prouhet--Tarry--Escott pairs). Twins
aligned at three shared positions under the row pairing are exactly
the perturbations of the square confined to two columns; twins aligned
at two are the skew exchanges. Empirically the two realizations
exclude each other for every tested pairing of every column-block
pair arising from the fifteen squares surveyed, across more than
$420{,}000$ explicit pairings. We prove (Theorem~\ref{thm:G}) that
the exclusion holds for \emph{every} pairing whenever the pair of
sextuples admits at most one twin up to complementation --- which
holds, with no exception, for all $1{,}078$ column-block pairs
arising from the fifteen squares surveyed. An exhaustive census of
\emph{free} bimagic pairs over $[1,N]$, $N\le26$, then identifies,
within the census range, the two residual forms at which the general
conjectures fail: pairs with four twins occur, in the two incidence
forms singled out by our case analysis; their alignment statistics
depend only on the form, as reproduced by an abstract six-point
model, and one form admits mixed pairings whereas the other has
none. Neither form occurs among the column blocks of the fifteen
squares surveyed. The uniqueness condition holds throughout the
surveyed square environment and is conjectured as an environmental
property.
\end{abstract}

\section{Introduction}

A magic square is \emph{bimagic} if it remains magic after squaring each
entry, and \emph{semi-bimagic} if rows and columns (but not necessarily
the diagonals) have constant sum and constant sum of squares. The first
semi-bimagic square of order six with distinct integers was published by
G.~Pfeffermann in 1894~\cite{pfeffermann}; it uses integers up to $49$,
with $S_1=150$ and $S_2=5150$, and has remained the reference artifact
for $132$ years. The systematic record of this field is maintained by
Christian Boyer~\cite{multimagie}.

For \emph{fully} bimagic squares of order six with distinct integers,
the minimal known example is due to Lee Morgenstern (2006):
$(\mathrm{MaxNb},S_1,S_2)=(72,219,10663)$, proved minimal in
2008--09~\cite{multimagie}. To our knowledge, the corresponding minimality
question for the \emph{semi}-bimagic case --- how small can the largest
entry be when the diagonal constraints are dropped --- has not been
recorded, either as a problem or as a record~\cite{multimagie}.

Our contribution is the answer at order six:
$\mathrm{MaxNb}=42$, with the full classification of the optimal sets.
The pair of results (Morgenstern's $(72,219,10663)$ versus our
$(42,129,3659)$) also quantifies, for the first time at order six, the
cost of the two diagonal constraints.

The second half of the paper (Section~\ref{sec:twins}) turns from the
record to the mechanism. The enumeration shows that survival is
decided almost entirely at the orthogonality stage, and
Remark~\ref{rem:richness} that no simple counting invariant governs
it. We isolate the local objects that do carry structure: \emph{twin
triples}, degree-$2$ Prouhet--Tarry--Escott pairs between two disjoint
bimagic sextuples, whose two possible realizations under a pairing
(row-aligned and skew) we prove to correspond bijectively to the
homogeneous and affine local exchanges of a square
(Proposition~\ref{prop:real}). On every pair arising from the
fifteen squares surveyed, the two realizations exclude each other.
Theorem~\ref{thm:G} proves this exclusion for every pairing whenever
the pair carries at most one twin up to complementation, and a
screen of all $1{,}078$ column-block pairs across those squares
finds no exception to that uniqueness. On free pairs of bimagic
sextuples the picture is sharper still: an exhaustive census up to
$[1,26]$ (Section~\ref{sec:free}) shows that four-twin pairs occur,
in exactly the two configurations left open by the case analysis,
and that within the census the general exclusion fails on one of
them and holds on the other --- uniqueness is thus conjectured as a
property of the \emph{square environment}
(Conjecture~\ref{conj:world}).

Exhaustive enumeration is the accepted standard of proof in this area:
the impossibility of a $7\times7$ bimagic square of consecutive integers
is accepted on the basis of four independent programs, and Morgenstern's
minimality on months of dedicated computation~\cite{multimagie}. We meet
that bar by running two algorithms with different assumptions on every
exclusion (Section~\ref{sec:method}).

\section{Definitions and statements}

\begin{definition}
A \emph{semi-bimagic square of order $6$} is a $6\times6$ array of
pairwise distinct positive integers such that all rows and all columns
have the same sum $S_1$ and the same sum of squares $S_2$. Its
$\mathrm{MaxNb}$ is its largest entry.
\end{definition}

\begin{theorem}[Minimality]\label{thm:min}
Semi-bimagic squares of order $6$ with $\mathrm{MaxNb}=42$ exist, and no
semi-bimagic square of order $6$ has $\mathrm{MaxNb}\le 41$.
\end{theorem}

\begin{theorem}[Classification at the optimum]\label{thm:class}
Every semi-bimagic square of order $6$ with $\mathrm{MaxNb}=42$ has
$S_1=129$, and its set of $36$ entries is one of exactly four sets,
namely $[1,42]$ minus one of
\[
\{19,\dots,24\},\quad \{4,11,18,25,32,39\},\quad
\{7,8,21,22,35,36\},\quad \{10,11,12,31,32,33\},
\]
with $S_2=3799,\ 3659,\ 3671,\ 3691$ respectively. All four sets are
closed under $x\mapsto 43-x$.
\end{theorem}

\begin{remark}[Independent replication]\label{rem:replication}
Theorem~\ref{thm:class} rests on two independent implementations. The
second (a different framing, written independently) enumerated all
candidate sets over the \emph{full} admissible window
$S_1\in[112,147]$ without assuming $1$ in the set, examining every
candidate admitting at least one row partition ($45{,}985$ entries in
the published data). Its success condition was pre-registered before
the run returned: exactly four sets with a square, all at $S_1=129$,
exactly the four sets above --- a square found anywhere else would have
falsified the theorem. The outcome: the four sets above, with row
partition counts $132, 50, 83, 59$ matching the first implementation,
and no square elsewhere. A pre-registered negative control behaved as
required: the centred set omitting $\{2,3,4,39,40,41\}$ (with
$S_2=3459$, \emph{below} the realized minimum) admits six row
partitions and no square. As a further internal cross-check, the
per-slot totals of examined candidates ($84{,}948$; $76{,}605$;
$76{,}605$; $84{,}948$ over the four slots, $323{,}106$ in all)
satisfy an exact mirror identity: the map $x\mapsto 42-x$ on the
omitted $6$-subsets of $\{1,\dots,41\}$ (the entry $42$ belongs to
every candidate set) sends stratum $S_1$ to stratum $259-S_1$ and
preserves the candidacy conditions, pairing the slots as
$(112\text{--}126,\,133\text{--}147)$ and
$(127\text{--}129,\,130\text{--}132)$ --- the examined counts pair
exactly, while the partitionable counts rightly do not, since no
entrywise bijection transports the semi-bimagic structure.
\end{remark}

\section{The squares}\label{sec:squares}

The minimal-$S_2$ square, $(42,129,3659)$, on the set omitting
$\{4,11,18,25,32,39\}$ (an arithmetic progression of step $7$):
\[
\begin{pmatrix}
 1 & 12 & 20 & 27 & 33 & 36\\
28 & 31 & 37 &  2 & 10 & 21\\
38 & 23 & 35 & 14 & 16 &  3\\
13 & 15 &  5 & 40 & 22 & 34\\
19 &  6 &  8 & 29 & 41 & 26\\
30 & 42 & 24 & 17 &  7 &  9
\end{pmatrix}
\]

A square on the set omitting the six central consecutive integers
$\{19,\dots,24\}$, $(42,129,3799)$:
\[
\begin{pmatrix}
 1 &  6 & 26 & 29 & 33 & 34\\
27 & 28 & 31 & 36 &  2 &  5\\
32 & 35 &  3 &  4 & 25 & 30\\
40 & 39 & 17 & 14 & 12 &  7\\
11 &  8 & 42 & 37 & 16 & 15\\
18 & 13 & 10 &  9 & 41 & 38
\end{pmatrix}
\]

A $26/28$ square (two magic, non-bimagic diagonals, matching
Pfeffermann's diagonal pattern but with $\mathrm{MaxNb}=42$):
\[
\begin{pmatrix}
 1 & 26 & 29 & 34 & 33 &  6\\
27 & 31 & 36 &  5 &  2 & 28\\
32 &  3 &  4 & 30 & 25 & 35\\
18 & 10 &  9 & 38 & 41 & 13\\
11 & 42 & 37 & 15 & 16 &  8\\
40 & 17 & 14 &  7 & 12 & 39
\end{pmatrix}
\]

All three were re-verified by an independent checker calibrated on the
four historical squares (see Data availability).

\section{Method: two independent algorithms}\label{sec:method}

\textbf{Algorithm 1 (set-sweep).} Enumerates candidate $36$-element
sets, then partitions into six bimagic sextets (equal $S_1$ and $S_2$),
then column extensions with pruning. It does \emph{not} assume
$1\in$ set: no translation lemma is used.

\textbf{Algorithm 2 (bucket).} Enumerates all sextets of $[1,M]$,
buckets them by $(S_1,S_2)$, and searches for six pairwise disjoint
sextets within a bucket; it assumes $1\in$ set and $M\in$ set via the
translation lemma.

The two implementations therefore differ in their hypotheses, not only
in their code. Every exclusion below carries both.

\begin{center}
\small
\begin{tabular}{@{}l>{\raggedright\arraybackslash}p{2.5cm}>{\raggedright\arraybackslash}p{4.1cm}>{\raggedright\arraybackslash}p{4.9cm}@{}}
\toprule
$\mathrm{MaxNb}$ & outcome & algorithms & logs\\
\midrule
36 & excluded & set-sweep + bucket + brute-force sextets & \texttt{blim3yxim}, \texttt{groups36.log}, \texttt{verify36.log}\\
37 & excluded & set-sweep + bucket & \texttt{blim3yxim}, \texttt{bn0vyvt60}\\
38 & excluded & set-sweep + bucket & \texttt{blim3yxim}, \texttt{bn0vyvt60}\\
39 & excluded & set-sweep + bucket & \texttt{b62iegrr5}, \texttt{groups39.log}\\
40 & excluded & bucket + set-sweep ($82{,}251$ sets) & \texttt{bjzvtqu5c}, \texttt{bz9px7zym}\\
41 & excluded & bucket + set-sweep ($658{,}008$ sets) & \texttt{bjzvtqu5c}, \texttt{bz9px7zym}\\
42 & attained + classified & two independent implementations, full window (Remark~\ref{rem:replication}) & \texttt{bjzvtqu5c}, \texttt{b5zmw01yv}, \texttt{b7x4hoknj}, \texttt{bmlofdsyf}, \texttt{replica-full-s1\_*.json}, \texttt{replica-full-verdetto.log}\\
\bottomrule
\end{tabular}
\end{center}

\section{The four optimal sets and the collapse}\label{sec:collapse}

At $\mathrm{MaxNb}=42$, $S_1=129$ is forced. The exhaustive phase-2 run
found $324$ partitions of candidate sets into six bimagic rows, of which
exactly $24$ extend to a full square, distributed over the four sets:

\begin{center}
\begin{tabular}{llll}
\toprule
omitted set & $S_2$ & partitions (extending/total) & structure\\
\midrule
$\{19,\dots,24\}$        & 3799 & 8 / 59  & six central consecutive\\
$\{4,11,18,25,32,39\}$   & 3659 & 8 / 132 & arithmetic progression, step 7\\
$\{7,8,21,22,35,36\}$    & 3671 & 4 / 50  & three adjacent pairs\\
$\{10,11,12,31,32,33\}$  & 3691 & 4 / 83  & two blocks of three\\
\bottomrule
\end{tabular}
\end{center}

Both totals ($324$ and $24$) were replicated by an independent
implementation with a different framing.

\medskip
\noindent\textbf{Where candidates die.} Within the symmetric class, the
counting chain has four declarable stages:
\[
546 \ \text{arithmetic candidates}
\ \xrightarrow{\ \text{sextet cover}\ } 546
\ \xrightarrow{\ \text{row partitions}\ } 382
\ \xrightarrow{\ \text{orthogonality}\ } 4 .
\]
No candidate dies at the sextet-cover stage; $164$ die for lack of any
row partition; and $378$ of the remaining $382$ ($99.0\%$) die at the
column-extension (orthogonality) stage. The bottleneck is orthogonality,
not arithmetic. A striking witness: the centred set omitting
$\{2,3,4,39,40,41\}$ passes all arithmetic conditions with
$S_2=3459$ --- \emph{below} the realized minimum --- and admits six row
partitions, yet no square exists on it.

\medskip
\noindent\textbf{Symmetry as output.} All four optimal sets are closed
under $x\mapsto 43-x$. This is not a filter of the search: $S_1=129$
only forces the mean $21.5$, and the enumeration is symmetry-agnostic.
The central symmetry emerges.

\medskip
\noindent\textbf{A finer instance.} The omitted set of the
$S_2$-minimal winner is exactly the residue class $4\pmod 7$ of
$\{1,\dots,42\}$. The mirror $x\mapsto 43-x$ sends class $r$ to class
$1-r\pmod 7$, and $r=4$ is the \emph{unique} self-mirror class. Its
five sister classes are admissible candidates in their own right
($S_1(r)=133-r$, $S_2(r)=3815-35r-r^2$, always integral), yet none of
them admits a single row partition: among the six whole-class
omissions, the unique symmetric one is the only one to reach the
partition stage --- and it wins.

\begin{proposition}[Lifting of the mirror symmetry]\label{prop:lift}
The involution $\sigma:x\mapsto 43-x$ sends bimagic sextets of a
mirror-closed entry set to bimagic sextets
($\textstyle\sum(43-x)=258-129=129$;
$\sum(43-x)^2=11094-2\cdot 43\cdot 129+\sum x^2=\sum x^2$), hence acts
on the row partitions of each optimal set and on orthogonal pairs of
partitions, preserving the intersection profile. Orbits have size one
or two, so an odd count of orthogonal pairs forces a $\sigma$-fixed
pair, whose square satisfies $43-M=M$ up to row and column
permutations (\emph{self-complementary}); a non-fixed pair is
exchanged with its $\sigma$-image, the two squares being complements
of one another up to permutations, composed with transposition when
$\sigma$ swaps the row and column roles. With the orthogonal-pair
counts $5,5,2,2$: the $S_2=3799$ set carries one self-complementary
square and two plainly complement-paired orbits; $3659$ one
self-complementary square, one plain and one transposed orbit; $3671$
one plain orbit; $3691$ one transposed orbit (all verified
explicitly). The mirror symmetry of the optimal sets thus lifts to the
optimal squares themselves.
\end{proposition}

\section{A corollary on squares of squares}

Squaring the $(42,129,3659)$ square cell-wise gives a $6\times6$
\emph{semi-magic square of $36$ distinct perfect squares} with constant
$3659$ and largest entry $42^2=1764$. (Magic squares of squares of order
$6$ with diagonals are known: Boyer 2005~\cite{boyer2005}; the point
here is only the size of the entries, inherited from the record.)

\section{The local structure of orthogonality: twin triples}
\label{sec:twins}

The enumeration of Section~\ref{sec:collapse} shows that survival is
decided almost entirely at the orthogonality stage ($99.0\%$ of the
row-partitionable symmetric candidates die there), and
Remark~\ref{rem:richness} that neither symmetry nor partition
abundance governs that stage. This section opens the structural study
of the mechanism itself. The natural local objects are pairs of
columns: two columns of a semi-bimagic square are two disjoint
six-element sets with the same sum and the same sum of squares,
interacting through the rows. We first develop the combinatorics of
such pairs in isolation, then return to the square.

\subsection{Twins, alignment, realizations}

\begin{definition}[bimagic pair, twins]\label{def:twin}
A \emph{bimagic pair} is a pair $(A,B)$ of disjoint six-element sets
of positive integers with $\sum A=\sum B$ and $\sum_{a\in A}a^2
=\sum_{b\in B}b^2$. A \emph{twin} of $(A,B)$ is a pair
$g=(T_A,T_B)$ of three-element subsets $T_A\subset A$,
$T_B\subset B$ with
\[
\textstyle\sum T_A=\sum T_B \quad\text{and}\quad
\sum_{x\in T_A}x^2=\sum_{y\in T_B}y^2,
\]
i.e.\ an instance of the Prouhet--Tarry--Escott problem of degree $2$
and size $3$~\cite{pte} inside the pair. We write $G(A,B)$ for the
set of twins of $(A,B)$, and
$\bar g=(A\setminus T_A,\,B\setminus T_B)$ for the \emph{complement}
of a twin.
\end{definition}

Complementation is well defined: $\sum(A\setminus T_A)=\sum A-\sum
T_A=\sum B-\sum T_B=\sum(B\setminus T_B)$, and likewise for squares.
Since $T_A\ne A\setminus T_A$, the map $g\mapsto\bar g$ is a
fixed-point-free involution of $G(A,B)$; in particular $|G(A,B)|$ is
always even.

\begin{lemma}[rigidity]\label{lem:rigid}
Two three-element subsets of the same six-element set sharing two
elements and having the same sum coincide. Consequently, if two
distinct twins share exactly two elements on some side, their
invariants $(\sum T,\sum T^{(2)})$ differ.
\end{lemma}

\begin{proof}
The third element is determined already by the linear sum.
\end{proof}

\begin{definition}[pairing, alignment]\label{def:align}
A \emph{pairing} of a bimagic pair is a bijection $\pi\colon A\to B$;
when $A,B$ are two columns of a square, the rows provide one. The
\emph{alignment} of a twin $g=(T_A,T_B)$ under $\pi$ is
\[
k(g,\pi)\;=\;|\pi(T_A)\cap T_B|\;\in\;\{0,1,2,3\},
\]
and we write $N_3(\pi)$, $N_2(\pi)$ for the numbers of twins of
$(A,B)$ with $k=3$, resp.\ $k=2$.
\end{definition}

Twins aligned at $3$ and at $2$ are precisely the two ways a twin can
be \emph{realized} on a square. Fix two columns $c_1,c_2$ of a
semi-bimagic square $M$; its entries being pairwise distinct, so are
in particular the twelve paired entries $p_j=M_{j,c_1}$,
$q_j=M_{j,c_2}$ ($j$ the row) --- the hypothesis used below to
exclude $|J|\in\{1,2\}$ and to recover rows uniquely. Write
$d_j=q_j-p_j$, $e_j=q_j^2-p_j^2$; the row pairing is
$\pi(p_j)=q_j$.

\begin{proposition}[realizations]\label{prop:real}
\emph{(i) (row-aligned twins are the moves confined to two columns.)}
For a set of rows $J$ with $0<|J|<6$, swapping the entries $p_j$ and
$q_j$ in the rows $j\in J$ yields again a semi-bimagic square if and
only if $\sum_{j\in J}d_j=0$ and $\sum_{j\in J}e_j=0$
(\emph{homogeneous system}). Every such $J$ has $|J|=3$, and
\[
J\;\longmapsto\;\bigl(\{p_j\colon j\in J\},\,\{q_j\colon j\in J\}\bigr)
\]
is a bijection from the proper solutions of the homogeneous system
onto the twins with $k(g,\pi)=3$.

\emph{(ii) (skew twins.)} Call \emph{affine system} the set of
triples $(j_1,j_2,K)$ with $j_1\ne j_2$, $K$ a two-element subset of
the remaining four rows, and
\[
q_{j_1}-p_{j_2}+\textstyle\sum_{j\in K}d_j=0,\qquad
q_{j_1}^2-p_{j_2}^2+\textstyle\sum_{j\in K}e_j=0 .
\]
Then
$(j_1,j_2,K)\mapsto\bigl(\{p_{j_2}\}\cup\{p_j\colon j\in K\},\,
\{q_{j_1}\}\cup\{q_j\colon j\in K\}\bigr)$
is a bijection from the affine system onto the twins with
$k(g,\pi)=2$.
\end{proposition}

\begin{proof}
(i) Each row keeps its two entries under the swap, so all row sums
and row square-sums are unchanged, as are the other columns; column
$c_1$ changes by $\sum_J d_j$ in sum and $\sum_J e_j$ in square-sum
(and $c_2$ by the negatives), whence the criterion. Sizes: the full
set $J=\{1,\dots,6\}$ solves the system trivially (the two columns
have equal sums and equal square-sums), so complements of solutions
are solutions; $|J|=1$ forces
$d_j=e_j=0$, i.e.\ $p_j=q_j$, excluded; if $J=\{i,j\}$ has size
$2$, then $d_i=-d_j$, and $e=d\,(p+q)$ turns $\sum e=0$ into
$p_i+q_i=p_j+q_j$; combined with $d_i=-d_j$ this gives $q_i=p_j$ and
$p_i=q_j$, repeated entries --- excluded. So proper solutions have
$|J|=3$. The map lands on twins: the sums of
$\{p_j\}_J$ and $\{q_j\}_J$ agree iff $\sum_J d=0$, likewise for
squares; and $\pi(\{p_j\}_J)=\{q_j\}_J$ gives $k=3$. It is injective
($J$ is the set of rows of $T_A$), and surjective: if $k(g,\pi)=3$
then $\pi(T_A)=T_B$, so with $J$ the rows of $T_A$ we get
$\sum_J d=\sum T_B-\sum T_A=0$, and likewise for $e$.

(ii) The two displayed equations say exactly that the two triples
have equal sums and equal square-sums. Alignment: $\pi(T_A)
=\{q_{j_2}\}\cup\{q_j\}_{j\in K}$, and since $j_1\ne j_2$ and
$j_1,j_2\notin K$, the intersection with
$T_B=\{q_{j_1}\}\cup\{q_j\}_{j\in K}$ is $\{q_j\}_{j\in K}$, so
$k=2$. Conversely, from a twin with $k=2$ one recovers uniquely: $K$
= the rows $j$ with $p_j\in T_A$ and $q_j\in T_B$ (exactly two of
them), $j_2$ = the row of the element of $T_A$ not paired into
$T_B$, $j_1$ = the row of the element of $T_B$ not hit by
$\pi(T_A)$; and $j_1\ne j_2$, since $j_1=j_2$ would give $k=3$.
\end{proof}

\begin{remark}[machine check]\label{rem:bijcheck}
Both bijections were verified computationally by complete enumeration
of both sides on $32{,}400$ pairings ($45$ column pairs from three
sample squares of different families $\times$ all $720$ bijections):
$1{,}008$ homogeneous solutions against $1{,}008$ twins at $k=3$ and
$9{,}072$ affine solutions against $9{,}072$ twins at $k=2$, with
multiset equality and injectivity holding on every single pairing.
\end{remark}

\subsection{The exclusion phenomenon}

On the actual squares the two realizations never coexist. On all
$1{,}200$ column pairs of the seven squares in our test harness (two
of the optimal squares, an auxiliary square with entries up to $44$
produced during our exploratory searches, Pfeffermann's and Boyer's
squares, Wroblewski's square with row sum $648$, and Morgenstern's
with row sum $330$), the
homogeneous system is solvable on $64$ pairs --- every solution with
$|J|=3$, as Proposition~\ref{prop:real} forces --- and \emph{no pair
solves both systems}. The real row pairing is not what protects:
permuting the pairing exhaustively --- a pre-registered deterministic
shuffle, $3{,}600$ pairings --- homogeneous solutions do appear where
the real pairing has none, yet still no pairing solves both. Neither
does any other pairing: on $45$ column pairs the full sweep of all
$720$ bijections each ($32{,}400$ pairings) finds $N_3(\pi)\,
N_2(\pi)=0$ throughout, and a wider screen ($388{,}800$ pairings:
the $S_2=3799$ optimal square, Pfeffermann's, and Wroblewski's
row-sum-$648$ square --- three different structural families --- over
all $36$ oriented orthogonal views, all $15$ column pairs and all
$720$ bijections) finds the same. This motivates the general conjecture:

\begin{conjecture}[exclusion, general form]\label{conj:excl}
For every bimagic pair $(A,B)$ and every pairing $\pi$,
\[
N_3(\pi)\cdot N_2(\pi)\;=\;0 .
\]
\end{conjecture}

Section~\ref{sec:free} refutes Conjecture~\ref{conj:excl} in this
generality. The theorem below identifies the hypothesis under which
its conclusion does hold, for every pairing.

\subsection{Exclusion from uniqueness: Theorem G}

\begin{lemma}[complementation preserves alignment]\label{lem:compl}
For every twin $g$ and pairing $\pi$:\ \ $k(\bar g,\pi)=k(g,\pi)$.
\end{lemma}

\begin{proof}
$\pi(A\setminus T_A)=B\setminus\pi(T_A)$, so
\[
|\pi(A\setminus T_A)\cap(B\setminus T_B)|
=6-|\pi(T_A)\cup T_B|
=6-\bigl(3+3-k(g,\pi)\bigr)=k(g,\pi). \qedhere
\]
\end{proof}

\begin{theorem}[exclusion from uniqueness]\label{thm:G}
Let $(A,B)$ be a bimagic pair with $G(A,B)=\emptyset$ or
$G(A,B)=\{g,\bar g\}$ (at most one twin up to complementation). Then
$N_3(\pi)\,N_2(\pi)=0$ for every pairing $\pi$.
\end{theorem}

\begin{proof}
If $G=\emptyset$ there is nothing to align. Let $G=\{g,\bar g\}$.
Alignment under a fixed $\pi$ assigns to each twin a single value
$k$, and by Lemma~\ref{lem:compl} the two twins carry the same value.
If $N_3(\pi)>0$, some twin has $k=3$, hence both do; $G$ is exhausted
by twins aligned at $3$, and $N_2(\pi)=0$.
\end{proof}

The proof is pure combinatorics of finite sets: the integer values
enter only through the definition of $G$. Theorem~\ref{thm:G}
concentrates all the arithmetic of Conjecture~\ref{conj:excl} into a
statement about twelve integers:

\begin{conjecture}[uniqueness of twins, general form]\label{conj:uniq}
Every bimagic pair admits at most one twin up to complementation:
$G(A,B)=\emptyset$ or $G(A,B)=\{g,\bar g\}$.
\end{conjecture}

\begin{corollary}
The uniqueness condition implies the exclusion property:
Conjecture~\ref{conj:uniq} implies Conjecture~\ref{conj:excl}.
\end{corollary}

Conjecture~\ref{conj:uniq} too is refuted in general by the free
census (Section~\ref{sec:free}): four-twin pairs exist. What the
fifteen-square data support exhaustively --- and what we conjecture
--- is its environment form:

\begin{conjecture}[uniqueness in the square environment]\label{conj:world}
For every pair $(A,B)$ of column blocks of the partitions entering a
semi-bimagic square of order $6$:\ \ $G(A,B)=\emptyset$ or
$G(A,B)=\{g,\bar g\}$.
\end{conjecture}

By Theorem~\ref{thm:G}, Conjecture~\ref{conj:world} yields the
exclusion on every pairing throughout the square environment.

\subsection{Evidence, and the residual case}\label{sec:residual}

\textbf{Direct census.} On the $45$ column pairs of the three sample
views of Remark~\ref{rem:bijcheck}: $|G|\in\{0,2\}$ with no other
value ($31$ empty, $14$ non-empty), and every non-empty $G$ is a
complementary pair $\{g,\bar g\}$. On the wide base: among all
$1{,}078$ distinct pairs of column blocks taken from the partitions
entering orthogonal pairs across fifteen squares (the four optimal
squares, the auxiliary square with entries up to $44$, and ten
historical or recorded squares: Pfeffermann's, Boyer's, two of
Morgenstern's, four of Wroblewski's and two of Rouanet's):
\[
|G|=0\ \text{on}\ 811\ \text{pairs},\qquad
|G|=2\ \text{on}\ 267\ \text{pairs},
\]
always a complementary pair, with no violation. The status of this
evidence is declared: exhaustive in the sample, not a proof.

\begin{remark}[the sweep is universal]\label{rem:sweep}
For a pair with $G=\{g,\bar g\}$ the distribution of alignments over
all $720$ pairings is forced by pure combinatorics, independently of
the values: the number of bijections of two six-element sets carrying
a fixed $3$-subset onto a fixed $3$-subset with intersection exactly
$k$ is $36,\,324,\,324,\,36$ for $k=3,2,1,0$. Hence every non-empty
pair shows the same sweep --- $36$ pairings with $N_3>0$, $324$ with
$N_2>0$, $360$ with neither --- exactly as observed on all $14$
non-empty pairs of the census. The same four numbers reappear as the
row marginals of the alignment tables of every four-twin pair of the
free census (Section~\ref{sec:free}).
\end{remark}

\textbf{The residual case.} What remains between the evidence and
Conjecture~\ref{conj:uniq} is the exclusion of two \emph{independent}
twins --- two twins that are neither equal nor complementary. The
route through complements, suggested to us in review, closes most of
the possible configurations:

\begin{proposition}[reduction of the residual case]\label{prop:reduction}
Let $g=(T_A,T_B)$ and $g'=(T'_A,T'_B)$ be \emph{independent} twins of
a bimagic pair $(A,B)$ --- twins that are neither equal nor
complementary, $g'\notin\{g,\bar g\}$ --- and
$(a,b)=(|T_A\cap T'_A|,\,|T_B\cap T'_B|)$.
Up to replacing $g'$ by $\bar g'$ (which turns $(a,b)$ into
$(3-a,3-b)$) and swapping the two sides, only two configurations are
possible:
\emph{(I) half-sum}, $(a,b)=(3,0)$: both triples of $g$ have halved
invariants, $2\sum T=\sum A$ and $2\sum T^{(2)}=\sum A^{(2)}$, and
then $|G(A,B)|\ge 4$;
\emph{(II) transversal}, $(a,b)=(2,1)$.
Consequently Conjecture~\ref{conj:uniq} holds if and only if no
bimagic pair carries a twin with halved invariants or a transversal
pair of twins.
\end{proposition}

\begin{proof}
The cells $(3,3)$ and $(0,0)$ are $g'=g$ and $g'=\bar g$, excluded by
independence. Let $a=3$, i.e.\ $T'_A=T_A$ (this covers $a=0$ by
complementation and $b\in\{0,3\}$ by the side swap): then $T'_B$ and
$T_B$ have the same invariants, both equal to those of $T_A$. Now
$b=2$ is impossible by Lemma~\ref{lem:rigid}. For $b=1$, write
$T_B=\{x,u,v\}$ and $T'_B=\{x,r,s\}$ with $x$ the shared element:
equal sums give $u+v=r+s$, equal square-sums give
$u^2+v^2=r^2+s^2$, hence $2uv=(u+v)^2-(u^2+v^2)=2rs$; so $\{u,v\}$
and $\{r,s\}$ are the root pairs of the same quadratic
$X^2-(u+v)X+uv$, i.e.\ $\{u,v\}=\{r,s\}$ and $T'_B=T_B$ ---
contradicting $b=1$. Finally $b=3$ gives $g'=g$. So $b=0$:
$T'_B=B\setminus T_B$, whose invariants equal those of $T_B$ exactly
when both are halved --- configuration (I) --- and then $g$, $\bar g$,
$(T_A,B\setminus T_B)$, $(A\setminus T_A,T_B)$ are four distinct
twins. Let $a,b\in\{1,2\}$: for $(a,b)=(2,2)$, subtracting the twin
conditions of $g'$ from those of $g$ the shared parts cancel, leaving
$a_3-a_4=b_3-b_4$ and $a_3^2-a_4^2=b_3^2-b_4^2$ for the third
elements, with $a_3\ne a_4$; hence $a_3+a_4=b_3+b_4$, so $a_3=b_3$
--- impossible, as $A\cap B=\emptyset$. The cell $(1,1)$ reduces to
$(2,2)$ by complementation, and $(2,1)\sim(1,2)$ is configuration
(II), where the same subtraction leaves a difference of single
elements against a difference of pairs and the collapse does not
bite.
\end{proof}

The census backs the case analysis cell by cell. On the $267$
non-empty pairs of the wide base: no twin with halved invariants, and
no independent second twin --- the environment uniqueness re-verified
through the case analysis itself. Moreover, the cells closed by the
linear argument alone are empty already at the linear level (as the
proof demands), while the lemma-closed cells and the two residual
cells host $2{,}512$ \emph{linear} candidates in all ($882$ in the
transversal cell, $194$ in the half-sum cell) --- every one killed by
the quadratic equation. The remaining cases are exactly
configurations (I) and (II) --- and the free census of the next
subsection (up to $N=26$) shows that \emph{both are realized} on
free bimagic pairs, with explicit members, and that in case (II)
there are pairings exhibiting both $k=3$ and $k=2$ alignment: the
two general conjectures are false. No exception occurs among the
$1{,}078$ column-block pairs surveyed; the case analysis identifies
the residual configurations that Conjecture~\ref{conj:world} asserts
the square environment excludes.

\begin{remark}[relation to the P.T.E. literature]
Twins are instances of the Prouhet--Tarry--Escott problem of degree
$2$ and size $3$, for which solutions abound in general. The point of
the uniqueness question is not the existence theory of P.T.E.\ pairs
but the way two \emph{bimagic sextuples} constrain the local P.T.E.\
structure: on free pairs the census finds the constraint stopping at
$|G|\le4$, in two rigid forms (Section~\ref{sec:free}), and
Conjecture~\ref{conj:world} says that inside squares it is total. We
are not aware of this question in the P.T.E.\ literature.
\end{remark}

\subsection{The free census: both residual configurations are
realized}\label{sec:free}

To test Conjectures~\ref{conj:excl} and~\ref{conj:uniq} outside the
square environment we ran an exhaustive census of \emph{free} bimagic
pairs: all pairs of disjoint six-element subsets of $[1,N]$ with
equal sum and equal sum of squares, for $N\le 26$ ($286{,}705$ pairs
at $N=26$).

\emph{Twin counts.} Throughout the census range $|G|\in\{0,2,4\}$:
four twins occur, never more. As forced by complementation, every
pair with $|G|=2$ consists of a complementary pair (verified
$11{,}158$ of $11{,}158$ at $N=24$). Four-twin pairs number $156$ in
$[1,22]$, $306$ in $[1,24]$, $564$ in $[1,26]$. Up to relabelling of
the four twins, they fall into the two incidence forms identified by
Proposition~\ref{prop:reduction}:
\emph{(II) transversal} ($156$, $300$, $552$ pairs), two complementary
twin pairs meeting transversally, and \emph{(I) half-sum} ($0$, $6$,
$12$ pairs, none below $[1,23]$), where both sextuples split into two
halves with halved invariants and the four twins are the $2\times2$
pairings of the halves.

\emph{The exclusion property fails for the transversal form, whose
realizations admit pairings with $N_3(\pi)>0$ and $N_2(\pi)>0$.}
Every transversal pair in the census shows the same alignment
statistics over its $720$ pairings: $32$ mixed ($N_3>0$ and
$N_2>0$), $40$ purely row-aligned, $488$ purely skew, $160$ null.
The smallest counterexample lies already in $[1,15]$:
$A=\{1,3,9,10,11,14\}$, $B=\{2,5,6,7,13,15\}$, whose
order-preserving pairing has $N_3=2$ and $N_2=2$. Every half-sum
pair in the census, by contrast, has alignment vector
$(k,3-k,3-k,k)$ over its four twins: $72$ pairings purely
row-aligned, $648$ purely skew, none mixed and none null --- the
exclusion holds for every pairing of a half-sum pair and fails for
the transversal form. These statistics depend only on the incidence
form, as shown by the abstract six-point model built from the
incidence matrix, which reproduces both tables by finite counting
over all realizations of each form, with the row marginals
$36/324/324/36$ of Remark~\ref{rem:sweep}.

\emph{The generator of the half-sum form.} Half-sum pairs arise from
quadruples of pairwise disjoint triples with equal sum and equal sum
of squares: each quadruple assembles into two sextuples in three
ways, giving three four-twin pairs, and within the census
\emph{every} half-sum pair arises this way. Quadruples number $0$ in
$[1,22]$, $1$ in $[1,23]$, then one more per step up to $4$ in
$[1,26]$ --- all four quadruples found in the census are translates
of the single seed
$\{1,17,18\},\{2,13,21\},\{3,11,22\},\{6,7,23\}$, four triples
sharing the invariants $(36,614)$.

\emph{The boundary.} Within the free census range both residual
configurations are realized and both general conjectures are
refuted; on the $1{,}078$ column-block pairs of the fifteen squares
surveyed, $|G|\le 2$ with no exception. What separates the surveyed
square environment from the free census is therefore not the
existence of the two residual forms --- both occur in the census ---
but their absence from the surveyed column blocks
(Conjecture~\ref{conj:world}).

\section{Open problems}

\begin{enumerate}
\item The global minimum of $S_1$ (and of $S_2$) over all semi-bimagic
squares of order $6$ with distinct integers, with $\mathrm{MaxNb}$
unconstrained. The trivial bound $S_1\ge 111$ improves for free:
$S_1=111$ forces the entry set $\{1,\dots,36\}$, whose maximum
$36\le 41$ is excluded by Theorem~\ref{thm:min}; hence $S_1\ge 112$.
\item The analogous minimality question at order $7$.
\item Prove Conjecture~\ref{conj:world} (uniqueness in the square
environment). By Theorem~\ref{thm:G} this establishes the exclusion
for every pairing throughout the environment. The two general
conjectures are refuted by the free census (Section~\ref{sec:free}),
which realizes both residual configurations of
Proposition~\ref{prop:reduction}; the open question is what
mechanism, if any, forces $|G|\le 2$ in the square environment.
\end{enumerate}

\begin{remark}[What the data already rule out]\label{rem:richness}
A natural probabilistic reading of the sieve --- \emph{the four
surviving sets are simply the richest in row partitions} --- fails.
Mirror symmetry is empirically necessary for the optimal sets, but far
from sufficient: $378$ of the $382$ mirror-symmetric candidates
admitting row partitions yield no square. Exhaustive enumeration
further shows that row-partition abundance alone cannot characterize
survival: the global champion among all $45{,}985$ candidates omits
$\{1,2,3,4,15,32\}$ (asymmetric, $S_1=141$) and admits $197$ row
partitions yet no square; $326$ non-winners are richer than the
poorest winner; and the four surviving sets admit between $50$ and
$132$ row partitions, of which only $4$--$8$ extend to a square.
Restricted to the mirror-symmetric candidates, the winners rank $1$,
$8$, $18$, $24$ by partition count. The unresolved problem is to
identify the global compatibility invariant governing the intersection
between row and column decompositions.
\end{remark}

\section*{Data availability}

All enumeration code, raw logs (with the identifiers cited in the tables
above), and the machine-readable list of the $324$ partitions are
available at
\url{https://github.com/Rekigan/semibimagic-6x6}; its README
maps every result of this paper to the script and the log that support
it. For Section~\ref{sec:twins}: \texttt{seduta-g.log}
(twin census, realization check, exhaustive $32{,}400$ sweep),
\texttt{selezione-j3.log} (the $1{,}200$-pair dichotomy),
\texttt{m5-predizioni.log} (pre-registered shuffle),
\texttt{terza-biunivocita.log} and \texttt{terza-vaglio-unicita.log}
(independent re-implementations of the bijection check and of the
$1{,}078$-pair uniqueness screen), and \texttt{residuo-ab.log} (the
case-analysis census of Proposition~\ref{prop:reduction}). For the
free census of Section~\ref{sec:free}: \texttt{fresco-vaglio.log} and
\texttt{fresco-excl.log} (census over $[1,22]$ and refutation sweep),
\texttt{invariante-156.log}, \texttt{invariante-n24.log} and
\texttt{salita-n26.log} (four-twin census to $N=26$: incidence forms
and alignment statistics), \texttt{halfsum-generatore.log} (the
quadruple generator), \texttt{modello-tabella-fine.log} (the
abstract-model alignment tables), and \texttt{cintura-finale.log}
(the complementation check at $|G|=2$).

\section*{Use of AI tools}

The computations and the writing of the code were done with AI tools,
under the direction and verification of the author.

\begin{thebibliography}{9}
\bibitem{pfeffermann} G.~Pfeffermann, Probl\`eme 1506: Carr\'e
magique de 6 \`a deux degr\'es (imparfait), \emph{Les Tablettes du
Chercheur}, Paris, 15 March 1894; solution, 15 April 1894.
\bibitem{multimagie} C.~Boyer, \emph{Multimagic squares site},
\texttt{multimagie.com} (accessed 2026-08-07): pages ``Unsolved
multimagic problems'', ``Smallest bimagic squares''.
\bibitem{boyer2005} C.~Boyer, Some notes on the magic squares of squares
problem, \emph{The Mathematical Intelligencer} \textbf{27} (2005).
\bibitem{pte} P.~Borwein and C.~Ingalls, The Prouhet--Tarry--Escott
problem revisited, \emph{Enseign. Math.} (2) \textbf{40} (1994),
3--27.
\end{thebibliography}

\end{document}
